\usepackage{graphicx} % Required for inserting images.
\usepackage[spanish,es-tabla]{babel} % Cambiar a español, para que diga tabla y no cuadro en las tablas.
\usepackage[utf8]{inputenc} % Para evitar problemas con signos, UTF-8 es estándar de codificación de caracteres más utilizado en la web.
\usepackage{parskip} % Para quitar la tabulación al inicio de cada parrafo.
\usepackage{float} % Para forzar el lugar de la imagen al insertarla usando [H].
\usepackage{hyperref} % Para poder agregar hipervinculos y los vinculos en el índice automáticamente.
% [colorlinks=ture] Cambia de color los diferentes tipos se vinculos.
\hypersetup{
    colorlinks=true, % Nos permite usar colores.
    linkcolor=black, % Ecuaciones, tablas, indices.
    urlcolor=blue, % Hipervinculos.
    citecolor=black, % Citas(Referencias).
    % pdftitle={PDF1} % Nombre del pdf (Para cambiarlo).
}

\usepackage[x11names]{xcolor} % Para poder usar más colores y ver documentación: https://www.overleaf.com/learn/latex/Using_colors_in_LaTeX
\usepackage{listings} % Para insertar código en el texto.